\chapter*{Abstract}
\addcontentsline{toc}{chapter}{Abstract}

This report develops, from first principles, a complete mathematical and
software framework for the distributed real-time estimation and control of
clusters of unmanned aerial vehicles (UAVs) comprising two dynamically
dissimilar vehicle classes: multirotor \emph{hexacopters}, modelled with full
unit-quaternion attitude dynamics, and \emph{fixed-wing} aircraft, modelled with
the standard six-degree-of-freedom rigid-body aerodynamics of Beard and McLain
instantiated on the Aerosonde airframe as a small-scale surrogate for a
medium-altitude long-endurance platform.

The two vehicle models are derived in full: every reference frame, kinematic
relation, force and moment, and Jacobian is stated explicitly, and every symbol
is defined at the point of use. For the hexacopter we derive the Hamilton
quaternion kinematics $\dot{\quat{q}}=\tfrac12\,\quat{q}\qmul[0,\vv{\omega}]$,
the Newton--Euler translational and rotational dynamics, the rotor
thrust/torque model, and the control-allocation mixer that maps six rotor
speeds to the collective-thrust/body-torque wrench, including its rank
structure and single-rotor-failure redundancy. For the fixed-wing aircraft we
derive the wind triangle, the sigmoid-blended stall aerodynamics, the full
force and moment coefficient expansions, and the coupled rotational dynamics in
the product-of-inertia $\Gamma$ form, together with a numerically computed
straight-and-level trim.

Each agent runs two constrained real-time-iteration solvers as its onboard
computational core: a nonlinear model-predictive \emph{controller} and a
moving-horizon \emph{estimator}, both of which impose general nonlinear state and
input constraints on the full problem and execute in the real-time-iteration
pattern, one iteration per sample, with the measurement-to-input path reduced to
a single affine evaluation. Around this core we construct
a distributed moving-horizon estimation (DMHE) layer, in which agents fuse
relative inter-agent measurements to remain localised through GNSS outages, and
a distributed model-predictive control (DMPC) layer that couples the agents
through formation objectives and control-barrier-function collision-avoidance
constraints, coordinated over a communication graph by consensus and
alternating-direction methods.

A dependency-free reference implementation---C for the numerical core, C++ for
the coordination layer, and Python for simulation---accompanies the report. In
closed-loop simulation on the full nonlinear plants, a six-hexacopter cluster
assembles and reconfigures a formation while maintaining collision-free
separation, and the DMHE keeps GNSS-denied agents localised to within
\SI{9}{\centi\meter} of ground truth through a ten-second blackout during which a
dead-reckoning filter drifts past \SI{1.2}{\meter}; a four-aircraft fixed-wing
cluster holds a coordinated-turn echelon formation. A $126$-run Monte-Carlo
campaign---over noise realisations, GNSS-denial severities from one to five of
six agents, degraded communication graphs and an estimator ablation---records
zero divergences and zero solver failures, a 95th-percentile denied-agent error
below \SI{0.12}{\meter} whenever every denied agent retains a localised
neighbour, and isolates \emph{anchor coverage} of the communication graph as the
condition on which that resilience rests. A computational study instantiates the
embedded solvers on the drone models themselves: one constrained
real-time-iteration cycle of the full quaternion-hexacopter NMPC executes in
about a hundred microseconds re-linearised, or below ten microseconds median
in its warm-started mode, with a \SI{30}{\nano\second}
measurement-to-input latency and accuracy identical across modes. The embedded
implementations are finally substituted into the closed loop itself: over
paired-seed cluster campaigns the deployed C estimator and controller
reproduce the reference behaviour with zero safety violations while
\emph{improving} the estimation error to six centimetres, and a leveled
multi-hop anchoring extension of the DMHE restores decimetre accuracy in the
anchor-coverage-broken cases, so the complete distributed
estimation-and-control loop is validated with the embedded solvers in it on
the nonlinear multi-agent plant. Every reported number is produced by the
accompanying code.
